\heading{高等数学A（II）-模拟题1-期中}

\noteinfo{题目和答案由AI生成。}

\begin{question}
计算积分
\[
I=\iint_D \sqrt{|y-x^2|}\,dx\,dy,
\qquad D=[-1,1]\times[-1,2].
\]

\begin{solution}

按 $y=x^2$ 分段。对固定 $x$，
\[
\int_{-1}^{x^2}\sqrt{x^2-y}\,dy+\int_{x^2}^{2}\sqrt{y-x^2}\,dy
=\frac23\bigl((x^2+1)^{3/2}+(2-x^2)^{3/2}\bigr).
\]
故
\[
I=\frac23\int_{-1}^{1}\bigl((x^2+1)^{3/2}+(2-x^2)^{3/2}\bigr)dx.
\]
利用偶性，
\[
I=\frac43\int_0^1(x^2+1)^{3/2}dx+\frac43\int_0^1(2-x^2)^{3/2}dx.
\]
分别计算可得
\ansmath{I=\frac12\ln(1+\sqrt2)+\frac{13\sqrt2}{6}+\pi.}
\end{solution}
\end{question}

\begin{question}
计算曲线积分
\[
I=\int_{\Gamma}\frac{(x-\tfrac12-y)\,dx+(x-\tfrac12+y)\,dy}{(x-\tfrac12)^2+y^2},
\]
其中 $\Gamma$ 从 $(0,-1)$ 沿抛物线 $y^2=1-x$ 到 $(0,1)$。
\begin{solution}

取参数
\[
 x=1-t^2,\qquad y=t,\qquad -1\le t\le1,
\qquad dx=-2t\,dt,\quad dy=dt.
\]
代入得
\[
I=\int_{-1}^{1}\frac{2t^3+2t^2-\tfrac12t+\tfrac12}{t^4+\tfrac14}\,dt.
\]
分拆并直接积分可得
\[
I=\frac\pi2+2\arctan 3.
\]
因此
\ansmath{I=\frac\pi2+2\arctan 3.}
\end{solution}
\end{question}

\begin{question}
计算曲面积分
\[
\iint_{\Sigma}2xy\,dy\,dz-y^2\,dz\,dx-x^2\,dx\,dy,
\]
其中 $\Sigma$ 为抛物面 $z=x^2+y^2$ 在 $0\le z\le2$ 之间的侧面，取外侧。
\begin{solution}

对应向量场
\[
\mathbf F=(2xy,-y^2,-x^2).
\]
散度
\[
\nabla\cdot\mathbf F=2y-2y+0=0.
\]
把侧面与顶盖圆盘 $z=2$ 补成闭曲面，由 Gauss 定理知总通量为 $0$，故侧面通量等于顶盖通量的相反数。顶盖外法向量向上，因此
\[
\iint_{z=2,x^2+y^2\le2}\mathbf F\cdot d\mathbf S=
\iint_{x^2+y^2\le2}(-x^2)\,dxdy.
\]
由对称性
\[
\iint_{x^2+y^2\le2}x^2\,dxdy=\frac12\iint_{x^2+y^2\le2}(x^2+y^2)\,dxdy
=\frac12\int_0^{2\pi}\!d\theta\int_0^{\sqrt2}r^3dr=\pi.
\]
故顶盖通量为 $-\pi$，于是侧面通量为
\ansmath{\pi.}
\end{solution}
\end{question}

\begin{question}
求曲线
\[
 x^3+y^3=xy
\]
所围成闭区域的面积。
\begin{solution}

令 $y=tx$，则
\[
 x^3(1+t^3)=tx^2 \Longrightarrow x=\frac{t}{1+t^3},\qquad y=\frac{t^2}{1+t^3},\qquad t\in[0,\infty).
\]
闭曲线位于第一象限。用参数曲线面积公式
\[
S=\frac12\int_0^{\infty}(x\,dy-y\,dx).
\]
直接代入并化简，得到
\[
 x\,dy-y\,dx=\frac{t^2}{(1+t^3)^2}\,dt.
\]
故
\[
 S=\frac12\int_0^{\infty}\frac{t^2}{(1+t^3)^2}\,dt.
\]
令 $u=t^3$，则
\[
S=\frac16\int_0^{\infty}\frac{du}{(1+u)^2}=\frac16.
\]
故
\ansmath{S=\frac16.}
\end{solution}
\end{question}

\begin{question}
设区域
\[
\Omega=\{(x,y,z):x^2+y^2\le z\le 2\}
\]
绕 $z$ 轴对称。求体积分
\[
V=\iiint_{\Omega}\frac{1}{1+x^2+y^2+z^2}\,dV.
\]

\begin{solution}

改用柱坐标 $x=r\cos\theta,y=r\sin\theta$，区域为
\[
0\le r\le\sqrt2,
\qquad r^2\le z\le2,
\qquad 0\le\theta\le2\pi.
\]
故
\[
V=2\pi\int_0^{\sqrt2}r\,dr\int_{r^2}^{2}\frac{dz}{1+r^2+z^2}.
\]
先对 $z$ 积分得
\[
\int_{r^2}^{2}\frac{dz}{1+r^2+z^2}=\frac{1}{\sqrt{1+r^2}}\left[\arctan\frac{z}{\sqrt{1+r^2}}\right]_{z=r^2}^{2}.
\]
故
\ansmath{V=2\pi\int_0^{\sqrt2}\frac{r}{\sqrt{1+r^2}}\left(\arctan\frac{2}{\sqrt{1+r^2}}-\arctan\frac{r^2}{\sqrt{1+r^2}}\right)dr.}
这已是整洁的单变量表达；若再令 $u=\sqrt{1+r^2}$ 可继续化简。
\end{solution}
\end{question}

\begin{question}
求函数
\[
 f(x,y,z)=x^2+y^2+z^2
\]
在约束
\[
 x+y+z=1,
\qquad x^2+y^2= z
\]
下的极值。
\begin{solution}

由 $x+y+z=1$ 与 $z=x^2+y^2$ 得
\[
 x+y+x^2+y^2=1.
\]
令 $s=x+y,\ p=x^2+y^2$，则约束是 $s+p=1$，且 $z=p$。目标函数
\[
 f=p+z^2=p+p^2.
\]
又由对任意实数 $x,y$，有
\[
\frac{s^2}{2}\le p\le s^2.
\]
结合 $s=1-p$ 得可行条件
\[
\frac{(1-p)^2}{2}\le p\le (1-p)^2.
\]
解得
\[
 3-2\sqrt2\le p\le \frac12.
\]
而 $f=p+p^2$ 在 $p\ge0$ 上单调增加，因此
\[
 f_{\min}=p+p^2\Big|_{p=3-2\sqrt2}=10-7\sqrt2,
\qquad
 f_{\max}=\frac12+\frac14=\frac34.
\]
对应点为：
\begin{enumerate}[label=（\arabic*）,leftmargin=2.8em,itemsep=0.2em]
\item 最大值在 $x=y=\tfrac12,\ z=\tfrac12$ 处取得；
\item 最小值在 $x,y$ 为二次方程 $t^2-(2\sqrt2-2)t+(3-2\sqrt2)/2=0$ 的两根时取得，等价地可取
\[
(x,y,z)=\left(\frac{\sqrt2-1}{2},\frac{\sqrt2-1}{2},3-2\sqrt2\right)
\]
不对，应检查；由 $p=(1-p)^2$ 的等号条件要求一元为零，故可行点为
\[
(x,y,z)=(\sqrt{3-2\sqrt2},0,3-2\sqrt2)
\]
及其对称变体，且再满足 $x+z=1$，因此
\[
(x,y,z)=(\sqrt2-1,0,3-2\sqrt2)
\]
及交换 $x,y$ 的点。
\end{enumerate}
故
\ansmath{f_{\max}=\frac34,\qquad f_{\min}=10-7\sqrt2.}
\end{solution}
\end{question}


